%=====================================================================
\chapter{The Sixteen-Week Schedule}\label{app:schedule}
%=====================================================================
\normalfigures

This appendix is the \textbf{source of record} for the teaching schedule. The
one-page \texttt{DBS\_Fall2026\_Lecture\_Schedule.pdf} distributed to students
is generated from it, and the slide pipeline parses the table below. Change
this first; regenerate the one-page version afterwards.

Twenty-three chapters are taught across sixteen weeks, so seven weeks carry
two chapters. The pairings are not arbitrary: in every case the two are either
one topic the printed book splits for length --- the two halves of the ER
model, the two halves of decomposition --- or a pair in which the second is
unusable without the first.

Two weeks are deliberately left alone. \textbf{Chapter~18, Normal Forms, has a
week to itself}, although it would pair naturally with Chapter~17. Its
laboratory breaks the running database and rebuilds it, and that does not fit
in half a week alongside anything else. \textbf{Chapter~21, Indexes}, likewise:
the measurement at a million rows takes most of a laboratory session to run
and the rest of it to argue about.

\section{The Schedule}

% MACHINE-READ. The slide builder takes the FIRST longtable in this file and
% requires: five columns; a bare integer week; chapters as comma-separated
% digits or N--M ranges; literal --- for an empty assessment cell; and no
% chapter appearing in two weeks.
\rowcolors{2}{white}{lightbg}
% column 3 holds the bold head "Chapters", which needs 1.8 cm at 10.95 pt
\begin{longtable}{@{}L{1.1cm}L{3.2cm}L{1.8cm}L{4.4cm}L{3.9cm}@{}}
\toprule
\rowcolor{primary!15}
\textbf{Wk} & \textbf{Lecture topic} & \textbf{Chapters} & \textbf{Laboratory} &
\textbf{Assessment} \\
\midrule
\endhead
1 & Why not files; the three-level architecture & 1, 2 & Break a file-based
system on purpose; find the three levels in SQLite & --- \\
2 & Who does what; your first database & 3, 4 & Install nothing; create a
database and measure what it cost & --- \\
3 & The entity-relationship model & 5 & Entities and attributes for the
university, on paper & --- \\
4 & Cardinality, participation, weak entities & 6 & The university ER diagram,
finished and argued over & --- \\
5 & The relational model; keys and constraints & 7, 8 & Orphan rows in a
correctly declared schema & \textbf{Quiz 1 (Ch 1--6)} \\
6 & Mapping a diagram to tables & 9 & Your diagram becomes eight tables, by
rule & \textbf{Assignment 1 set} \\
7 & Creating and changing a schema & 10, 11 & \texttt{CREATE TABLE}; then
change one that is already full & --- \\
8 & Relational algebra & 12 & Midterm, then algebra by hand on small relations
& \textbf{Midterm (Ch 1--11)} \\
9 & Joins; selecting rows and NULL & 13, 14 & Lose 73 sections and 19{,}992
rows to joins and NULL, then find them all & \textbf{Assignment 1 due} \\
10 & Sub-queries; grouping and aggregation & 15, 16 & One question seven ways;
an average moved by a third & \textbf{Project proposal} \\
11 & Functional dependencies & 17 & The wide report table: its dependencies,
its closures and its key & --- \\
12 & Normal forms: 1NF to BCNF & 18 & Break the running database, then
normalise it and measure both & \textbf{Quiz 2 (Ch 12--17)} \\
13 & Decomposition; views and triggers & 19, 20 & A lossy split that invents
facts; a summary table caught lying & \textbf{Assignment 2 set} \\
14 & Indexes, and where they stop helping & 21 & One million rows, with and
without an index & --- \\
15 & Transactions, concurrency and recovery & 22 & Two writers, one database,
and a lost update & \textbf{Assignment 2 due} \\
16 & Security, NoSQL and the project & 23 & Project presentations &
\textbf{Project due; presentations} \\
\bottomrule
\end{longtable}

\section{How the Assessment Adds Up}

HEC publishes the assessment instruments for this course --- sessional
examination, home assignments, quizzes, project, presentations, final
examination --- without weights. The weights below are the department's, and
are stated here so that a student can see the whole shape of the semester on
one page.

\rowcolors{2}{white}{lightbg}
\begin{longtable}{@{}L{4.6cm}L{2.0cm}L{8.4cm}@{}}
\toprule
\rowcolor{primary!15}
\textbf{Instrument} & \textbf{Weight} & \textbf{What it examines} \\
\midrule
\endhead
Two quizzes & 10\% & Weeks 5 and 12. Short, closed-book, on the chapters named
in the schedule \\
Two assignments & 10\% & Set in weeks 6 and 13. Both are done on your own
database and submitted as a change to it plus a short write-up \\
\textbf{Laboratory} & \textbf{15\%} & Assessed weekly from Week~2. This course
carries a \textbf{credited laboratory} --- one of its four credit hours --- and
the marks follow the credit \\
Midterm examination & 20\% & Week 8, Chapters 1--11 \\
Semester project & 20\% & The database itself, built across the whole semester
(\chref{project}), submitted in week 16 with its presentation \\
Final examination & 25\% & All chapters, with emphasis after the midterm \\
\bottomrule
\end{longtable}

\begin{alertbox}[The laboratory is a quarter of the course, not a formality]
This is the only course in the series with a \emph{credited} laboratory: HEC
specifies 4 credit hours as 3+1 in 2017 and 2025, and as 4 (3--3) in 2023. One
of those four hours is the lab, and 15\% of the marks follow it.

\smallskip
That is also why the laboratories in this handout are longer than those in its
sibling courses. They are not homework to be fitted around the lectures; they
are a timetabled hour with a demonstrator in the room, and every one of them
adds something to the database you will submit in Week~16.

\smallskip
\chref{project} states what is required and how it is marked. Read it in
Week~3, not in Week~15.
\end{alertbox}
